\documentclass{article}
\usepackage[utf8]{inputenc}
\usepackage[T2A]{fontenc}
\usepackage{helvet}
\usepackage{geometry}
\usepackage{xcolor}
\usepackage{eso-pic}
\usepackage{fancyhdr}
\usepackage{parskip}
\usepackage{microtype}
\usepackage[english, ukrainian]{babel}
\usepackage{url}
\usepackage{amsmath}
\usepackage{amssymb}
\usepackage{amsthm}
\usepackage{longtable}
\usepackage{booktabs}
\usepackage{array}
\usepackage{hyperref}
\usepackage{titlesec}

% Define brand colors
\definecolor{primary}{HTML}{293757}
\definecolor{footergray}{gray}{0.65}

\geometry{a4paper,left=3cm,right=3cm,top=3cm,bottom=3cm}
\renewcommand{\familydefault}{\sfdefault}
\setcounter{secnumdepth}{3}

% Font shapes for T2A/Helvetica (mapping to cmss for Cyrillic support)
\DeclareFontFamily{T2A}{phv}{}
\DeclareFontShape{T2A}{phv}{m}{n}{<->ssub * cmss/m/n}{}
\DeclareFontShape{T2A}{phv}{bx}{n}{<->ssub * cmss/bx/n}{}
\DeclareFontShape{T2A}{phv}{m}{it}{<->ssub * cmss/m/it}{}
\DeclareFontShape{T2A}{phv}{bx}{it}{<->ssub * cmss/bx/it}{}

\titleformat{\section}
  {\color{primary}\Huge\bfseries}{\thesection.}{1em}{}
\titlespacing*{\section}{0pt}{30pt}{12pt}
\titleformat{\subsection}
  {\color{primary}\LARGE\bfseries}{\thesubsection.}{1em}{}
\titlespacing*{\subsection}{0pt}{20pt}{8pt}
\titleformat{\subsubsection}
  {\color{primary}\Large\bfseries}{\thesubsubsection.}{1em}{}
\titlespacing*{\subsubsection}{0pt}{10pt}{6pt}

\fancyhf{}
\fancyhead[R]{\small\color{footergray} 29 червня 2026}
\fancyfoot[L]{\small\color{footergray} Технічні вимоги ERP/1. Модуль Аналітика}
\fancyfoot[R]{\small\color{footergray}\thepage}

\newcommand\HUGE[1]{\fontsize{38}{38}\selectfont #1}
\renewcommand{\headrulewidth}{0pt}
\renewcommand{\footrulewidth}{0.4pt}
\renewcommand{\footrule}{\color{footergray}\hrule width \textwidth height 0.4pt}

\begin{document}
\pagestyle{fancy}

\begin{titlepage}
\AddToShipoutPictureBG*{\AtPageLowerLeft{\color{primary}\rule{\paperwidth}{\paperheight}}}
\color{white}\thispagestyle{empty}\vspace*{4cm}\centering
  {\Huge \textbf{ERP/1: Аналітика} \par} \vspace{2cm}
  {\HUGE \textbf{Технічні вимоги} \par} \vspace{2.5cm}
  {\LARGE \textbf{Локальна DWH-інфраструктура \\ для аналізу та обробки великих даних} \par} \vspace{2.5cm}
\vfill
  {\large 2026 \copyright\ Системи Електронного Урядування \par}
\end{titlepage}

\newpage
\begin{center}
  {\Huge\color{primary}\bfseries Зміст\par}
\end{center}
\vspace{40pt}
\tableofcontents

\newpage
\begin{abstract}
У цьому документі наведено детальні технічні вимоги до підсистеми ERP/1: Аналітика, яка розробляється як сучасний високопродуктивний сервіс локального сховища великих даних (DWH), аналітичних обчислень та побудови бізнес-звітів (BI). Платформа інтегрується з децентралізованою екосистемою ERP/1 та передбачає локальне виконання швидких OLAP-запитів без передачі конфіденційної інформації за межі установ. Архітектурні вимоги фокусуються на використанні C99-сумісних мінімалістичних технологій без застосування екосистеми Rust, з повною інтеграцією в Erlang/OTP середовище периферійних вузлів. Документ визначає апаратний стек, вимоги безпеки відповідно до КСЗІ та NIST SP 800-53, а також містить порівняльний аналіз цільових інструментів для вибору рішення.
\end{abstract}

\newpage
\section{Умовні скорочення та визначення}

Терміни та скорочення, що використовуються в цьому документі:

\begin{longtable}{p{4cm}p{10cm}}
\toprule
\textbf{Термін / Скорочення} & \textbf{Значення} \\
\midrule
DWH & Data Warehouse (Сховище даних) \\
OLAP & Online Analytical Processing (Аналітична обробка в реальному часі) \\
ETL & Extract, Transform, Load (Витягнення, перетворення, завантаження) \\
ELT & Extract, Load, Transform (Витягнення, завантаження, перетворення) \\
BI & Business Intelligence (Бізнес-аналітика та візуалізація) \\
NIF & Native Implemented Function (Вбудована Erlang-функція мовою C) \\
OTP & Open Telecom Platform (Програмна екосистема Erlang) \\
Parquet & Відкритий стовпчиковий формат зберігання даних \\
AVX-512 & Розширення векторних інструкцій процесорів x86-64 \\
КСЗІ & Комплексна система захисту інформації \\
\bottomrule
\end{longtable}

\newpage
\section{Загальні відомості}

\subsection{Передумови}
Аналіз та побудова звітів за великими масивами транзакційних даних, реєстрів та документів вимагає виділення окремого обчислювального контуру, оптимізованого під читання великої кількості записів. Водночас, обробка даних у державному та корпоративному секторах накладає такі суворі архітектурні обмеження:
\begin{enumerate}
    \item \textbf{Конфіденційність та автономність}: Відповідно до вимог законодавства України щодо захисту персональних даних та державних реєстрів, аналітична обробка повинна відбуватися виключно в межах локального захищеного контуру (on-premises) установи без надсилання даних у хмари.
    \item \textbf{Мінімізація залежностей стеку компиляції}: Для спрощення сертифікації за вимогами КСЗІ та розгортання на великій кількості периферійних пристроїв необхідно обмежити використання складних інструментаріїв збирання та компіляції (виключити Rust/Cargo), віддаючи перевагу перевіреним системним засобам C99.
    \item \textbf{Інтеграція з Erlang/OTP}: - Аналітична підсистема повинна безшовно взаємодіяти з транзакційним контуром ERP/1, написаним на Erlang/Elixir, через нативні механізми портів або C-інтерфейси NIF.
\end{enumerate}

\subsection{Питання, що вирішуються}
Впровадження модуля «Аналітика» вирішує такі ключові завдання:
\begin{itemize}
    \item Усунення аналітичного навантаження (великих агрегаційних запитів) з транзакційної бази даних Mnesia/KVS.
    \item Швидкий аналіз історичних даних обсягом до десятків мільярдів записів безпосередньо на периферійних вузлах.
    \item Зберігання та обробка даних у стиснутому стовпчиковому форматі, що зменшує вимоги до дискового простору.
    \item Візуалізація звітів за допомогою інтегрованих інструментів бізнес-аналітики.
\end{itemize}

\section{Призначення та цілі впровадження}
Основною метою створення модуля є надання інструментів побудови аналітичних звітів, агрегації показників роботи установи (наприклад, статистика судових справ, фінансовий облік, обробка звернень) у реальному часі на локальних вузлах з використанням сумісних, легкозамінних та швидкодіючих відкритих технологій.

\newpage
\section{Класифікація вимог}

\subsection{Функціональні вимоги}

\subsubsection{Вимоги до локального аналітичного сховища (DWH) та його архітектури}
Система повинна реалізовувати локальне стовпчикове сховище даних на базі вбудовуваного аналітичного двигуна (DuckDB), що відповідає таким вимогам:
\begin{itemize}
    \item \textbf{Стовпчикове зберігання та стиснення}: зберігання агрегованих таблиць та журналів подій у стовпчиковому форматі для максимізації швидкості читання із застосуванням алгоритмів стиснення (ZSTD, LZ4).
    \item \textbf{Сумісність форматів}: повна підтримка роботи з відкритими аналітичними форматами файлів (Parquet, CSV, JSON).
    \item \textbf{Самостійність та ізоляція}: функціонування у вигляді окремого самостійного процесу (або ізольованого системного демона) з власним пулом виділеної оперативної пам'яті та лімітами обчислювальних ядер CPU, що запобігає впливу аналітичних обчислень на основні транзакційні процеси.
    \item \textbf{Масштабованість}: підтримка горизонтального масштабування шляхом федеративного виконання SQL-запитів над множиною незалежних баз та розподілених Parquet-файлів, а також реалізація механізмів Peer-to-Peer реплікації даних між периферійними вузлами для забезпечення відмовостійкості.
    \item \textbf{Автономність}: сховище повинно повноцінно функціонувати в Air-Gapped режимі без підключення до зовнішніх мережевих сервісів.
\end{itemize}

\subsubsection{Вимоги до інтеграції та ETL/ELT контурів}
Організація імпорту даних із транзакційної підсистеми ERP/1:
\begin{itemize}
    \item Асинхронний експорт змін транзакційних таблиць Mnesia/KVS у аналітичне сховище.
    \item Батчеве завантаження (Batch Loading) історичних масивів даних у фоновому режимі з низьким пріоритетом використання процесора (nice).
    \item Можливість виконання легких ELT перетворень безпосередньо всередині сховища за допомогою декларативних SQL-запитів.
\end{itemize}

\subsubsection{Вимоги до векторизованого виконання запитів та обчислень}
Аналітичний двигун повинен мати такі можливості:
\begin{itemize}
    \item Векторизоване виконання запитів (обробка даних пакетами рядків, а не по одному), що дозволяє повністю задіяти кеш процесора.
    \item Використання багатопотокового паралелізму для обробки одного аналітичного SQL-запиту.
    \item Наявність C99-сумісного API інтерфейсу для низькорівневої інтеграції з кодом Erlang/Elixir.
\end{itemize}

\subsection{Нефункціональні вимоги}

\subsubsection{Апаратні вимоги до обчислювальних вузлів}
В якості апаратної платформи використовуються стандартні сервери або робочі станції канцелярій:
\begin{itemize}
    \item Процесор: архітектура x86-64 з обов'язковою підтримкою векторних інструкцій AVX2 або AVX-512.
    \item Оперативна пам'ять: від 16 до 64 ГБ RAM (з оптимізацією під in-memory обчислення).
    \item Накопичувачі: NVMe SSD з високою швидкістю випадкового читання.
\end{itemize}

\subsubsection{Вимоги до програмного стеку}
Вимоги до технологічного стеку та системних обмежень:
\begin{itemize}
    \item \textbf{Використання C99}: Усі низькорівневі бібліотеки, драйвери, Erlang NIF модулі та інтерфейсні обгортки повинні бути розроблені на чистому стандарті мови C99 без складних зовнішніх бібліотек C++.
    \item \textbf{Erlang/OTP екосистема}: Керування життєвим циклом DWH, оркестрація ETL-процесів та обробка аналітичних задач повинні виконуватися як стандартні OTP-додатки (через супервізори та gen\_servers).
\end{itemize}

\subsubsection{Вимоги до продуктивності та масштабування}
Система аналітики повинна забезпечувати такі метрики продуктивності:
\begin{itemize}
    \item Швидкість сканування даних: не менше 100 мільйонів рядків на секунду на одному обчислювальному ядрі CPU.
    \item Час виконання складних агрегацій (GROUP BY, JOIN) на масивах до 100 млн записів: не більше 1.5 секунди.
    \item Можливість горизонтального масштабування шляхом об'єднання локальних периферійних баз у централізований аналітичний кластер.
\end{itemize}

\subsubsection{Вимоги до безпеки та захисту інформації (NIST SP 800-53)}
Захист аналітичних даних забезпечується такими засобами:
\begin{itemize}
    \item \textbf{SC-28 (Protection of Information at Rest)}: повне шифрування дисків за допомогою LUKS та інтеграції з TPM 2.0.
    \item \textbf{AC-3 (Access Enforcement)}: ізоляція доступу на рівні рядків таблиць (Row-Level Security) та атрибутивний контроль прав (ABAC).
    \item \textbf{Network Isolation}: блокування зовнішнього трафіку за допомогою Network Policies у Cilium.
\end{itemize}

\subsubsection{Вимоги до відмовостійкості та резервування}
Забезпечення доступності сервісу аналітики:
\begin{itemize}
    \item Підтримка механізму холодної та гарячої заміни вузлів без втрати історичних архівів.
    \item Асинхронне нічне резервне копіювання локальних даних у централізоване об'єктне сховище.
\end{itemize}

\newpage
\section{Додаток А. Порівняльний аналіз OLAP-платформ та засобів візуалізації}

Для вибору цільового рішення нижче наведено порівняльний аналіз двох провідних відкритих аналітичних баз даних (ClickHouse та MonetDB як найлегшого конкурента) та системи візуалізації звітів (Apache Superset) під кутом архітектурних вимог ERP/1:

\begin{longtable}{p{3cm} p{3.5cm} p{3.5cm} p{4cm}}
\toprule
\textbf{Критерій} & \textbf{ClickHouse} & \textbf{MonetDB} & \textbf{Apache Superset} \\
\midrule
\textbf{Тип системи} & Розподілена стовпчикова СКБД & Стовпчикова реляційна СКБД (standalone) & BI-платформа візуалізації та побудови дашбордів \\
\midrule
\textbf{Мова реалізації} & C++20 & Чистий C (C99) & Python (SQLAlchemy) + React (TypeScript) \\
\midrule
\textbf{Складність розгортання} & Висока (вимагає серверного кластера, Zookeeper/Keeper) & Низька (легковагий демон без зовнішніх залежностей) & Середня (запускається як окремий контейнер у K8s) \\
\midrule
\textbf{Інтеграція з Erlang/OTP} & Через TCP/HTTP клієнти або зовнішні драйвери & Через нативний Erlang-драйвер або C-порти & Через стандартні SQL-драйвери до ClickHouse/MonetDB \\
\midrule
\textbf{Цільове призначення} & Централізований аналітичний хаб для мільярдів рядків & Автономне легке аналітичне сховище для локальних вузлів & Клієнтський інтерфейс бізнес-аналітики \\
\bottomrule
\end{longtable}

\subsection{Рекомендації щодо архітектурного вибору}
\begin{enumerate}
    \item \textbf{Для децентралізованих периферійних вузлів (робочих станцій)} рекомендується використання \textbf{DuckDB}. Проте, згідно з технічними вимогами, рішення на базі DuckDB має бути повністю \textbf{самостійним} (ізольовані системні ресурси, окремі пули потоків виконання NIF/Port) та \textbf{масштабованим} (з підтримкою розподіленого федеративного виконання запитів і Peer-to-Peer синхронізації Parquet-файлів), що дозволяє успішно конкурувати з важкими клієнт-серверними рішеннями на локальному рівні.
    \item \textbf{Для центрального аналітичного вузла системи ERP/1} рекомендується розгортання кластера \textbf{ClickHouse}. Він бере на себе роль глобального DWH, куди стікаються агреговані дані з усіх 600+ установок.
    \item \textbf{Для побудови та відображення звітів користувачам} як окремий аналітичний додаток пропонується використовувати \textbf{Apache Superset}, який підключається як до локальних баз DuckDB, так і до центрального ClickHouse.
\end{enumerate}

\end{document}
